\documentclass[11pt]{article}
\usepackage{preamble}
\renewcommand{\answer}[1]{}

\begin{document}

\ladderheading{AP Statistics \textperiodcentered\ Topic 6.8 (CED 3.10) \textperiodcentered\ B: 2026-12-11 \textperiodcentered\ E: 2027-01-22}{One Difference, Not Two Separate Intervals}

\focusbanner{TODAY'S PROBLEM}

Town Oak has 2{,}200 households and Pine 1{,}800. Independent SRSs find 77 of 140 Oak households and 48 of 120 Pine households use curbside composting. Let $p_O-p_P$ mean Oak minus Pine.\par\smallskip \textbf{Elena:} ``The two-sample $z$-interval is $(0.030,0.270)$.''\par \textbf{Devon:} ``Subtract the farthest endpoints of separate 95\% intervals. That gives $(-0.020,0.320)$, so zero is plausible.''\par
\begin{center}
\textbf{Household samples}\par\smallskip
\begin{tabular}{|l|c|c|c|c|}
\hline
\textbf{Town} & \textbf{Use} & \textbf{Other} & \textbf{n} & \textbf{N} \\ \hline
\textbf{Oak} & 77 & 63 & 140 & 2200 \\ \hline
\textbf{Pine} & 48 & 72 & 120 & 1800 \\ \hline
\end{tabular}
\end{center}
\begin{firsttakebox}{FIRST TAKE --- before the video (5 min)}
Which method gives the requested 95\% interval for $p_O-p_P$? Cite one design or numerical detail.
\workspace{0.9in}
\end{firsttakebox}

\begin{callout}{\IconBook\ BUILD ONE INTERVAL FOR THE DIFFERENCE (keep on the page)}{calloutgreen}
For two independent random samples, estimate $p_1-p_2$ with a two-sample $z$-interval:
$(\hat p_1-\hat p_2)\pm z^*\sqrt{\hat p_1(1-\hat p_1)/n_1+\hat p_2(1-\hat p_2)/n_2}$.
Verify independent random samples or random assignment, $n_i\leq0.10N_i$ when sampling without replacement,
and at least 10 observed successes and failures in each group. For an interval, use each group's sample
proportion separately; because the samples are independent, estimated variances add before taking the square root.
Two separately constructed $C\%$ intervals do not automatically combine into one $C\%$ interval for $p_1-p_2$.
\end{callout}

\newpage
\textbf{Finish after the video:}\par\smallskip

\dokbadge{1}~\textbf{(a)} Define $p_O$ and $p_P$ and name the interval procedure for $p_O-p_P$.\par
\workspace{0.7in}

\dokbadge{2}~\textbf{(b)} Check independent samples, both 10 percent conditions, and all four counts. Compute the estimate, SE, margin of error, and 95\% interval.\par
\workspace{1.4in}

\dokbadge{3}~\textbf{(c)} Choose the warranted method. Use design, counts, variance structure, and zero's location. Explain why separate 95\% intervals do not produce the requested coverage and how Devon's result misleads.\par
\workspace{2.0in}

\begin{sentenceframebox}\raggedright\textbf{Frame:} The warranted method is \blank[1.0in] because the data come from \blank[2.2in], and the four counts are \blank[1.5in].\end{sentenceframebox}

\begin{sentenceframebox}\raggedright\textbf{Frame:} Devon's interval misleads a reader because \blank[2.4in]; as a result, zero appears \blank[0.8in] even though \blank[2.2in].\end{sentenceframebox}

\begin{summaryexitbox}{EXIT --- turn this in}
One thing the video changed about my first take: \hrulefill\par
\workspace{0.4in}
\end{summaryexitbox}

\end{document}
